% TOP_PROBLEM: {"id": 10000088, "title": "Uniqueness at the percolation uniqueness threshold", "category_id": 19, "category": {"id": 19, "name": "probability", "display_name": "Probability", "description": "Problems involving probability theory, stochastic processes, and random structures.", "slug": "probability", "order_index": 19, "created_at": "2026-07-31T15:26:25.671Z"}, "status": "open", "problem_number": "AMR-099-0088", "set_id": 13, "statusQualification": "Retains only the endpoint classification requested in the catalogue; the original paper uses site percolation, and its update resolves the strictly-above-threshold part."}
% FIRST_POSTED: 2026-10-06
% ATTEMPT_STATUS: unresolved
% OPENAI_CLAIM: {"schema_version":1,"source_repo":"https://github.com/openai/math","source_commit":"adc7f1241b42e322a6451854ab7e4b4c146bf78a","release_date":"2026-10-06","match":"partial","resolution":"partial","independently_reviewed":false,"second_pass":"n/a","families":[{"family":"213","title":"Critical percolation on every quasi-transitive graph","subject":"Probability and statistical mechanics","manuscripts":[{"title":"Critical bond and site percolation on the cubic lattice","pdf_path":"preprints/Critical-bond-and-site-percolation-on-the-cubic-lattice-September-24-2026/paper.pdf","readme_path":"preprints/Critical-bond-and-site-percolation-on-the-cubic-lattice-September-24-2026/README.md","date":"2026-09-24","theorem_ref":"thm:main (introduction.tex:17)","statement_sources":[{"path":"preprints/Critical-bond-and-site-percolation-on-the-cubic-lattice-September-24-2026/build/introduction.tex","line":17,"label":"thm:main"}]},{"title":"No percolation at criticality on quasi-transitive graphs","pdf_path":"preprints/No-percolation-at-criticality-on-quasi-transitive-graphs-September-24-2026/paper.pdf","readme_path":"preprints/No-percolation-at-criticality-on-quasi-transitive-graphs-September-24-2026/README.md","date":"2026-09-24","theorem_ref":"thm:main (sections/introduction.tex:26)","statement_sources":[{"path":"preprints/No-percolation-at-criticality-on-quasi-transitive-graphs-September-24-2026/build/sections/introduction.tex","line":26,"label":"thm:main"}]}],"lean":{"doc_path":"lean/docs/213.md","comparators":["lean/ComparatorChallenges/CriticalPercolation.json","lean/ComparatorChallenges/CriticalPercolation.lean","lean/ComparatorChallenges/CriticalZ3.json","lean/ComparatorChallenges/CriticalZ3.lean"],"declarations":["OAI.CriticalZ3.critical_no_infinite"],"covers_main_theorem":true},"reasoning_trace":null}]}
\documentclass[11pt]{article}
\usepackage[margin=25mm]{geometry}
\usepackage{fontspec}
\setmainfont{Latin Modern Roman}
\usepackage{amsmath,amssymb,mathtools}
\usepackage{unicode-math}
\setmathfont{Latin Modern Math}
\usepackage{enumitem}
\usepackage{xurl,hyperref,bookmark}
\hypersetup{colorlinks=true,urlcolor=blue}
\setcounter{secnumdepth}{0}
\setlength{\parindent}{0pt}
\setlength{\parskip}{5pt}
\setlength{\emergencystretch}{3em}
\begin{document}
\section*{Uniqueness at the percolation uniqueness threshold}

\subsection{OpenAI's claim}
The linked manuscripts present the following claims by OpenAI: Both bond and site percolation on $\mathbb Z^3$ are claimed to have only finite clusters at their respective critical probabilities; Critical bond percolation is claimed to have no infinite component on every infinite locally finite quasi-transitive graph with $p_c<1$. Together they provide the following support within the scopes described below.

\subsection{Mathematical statement of the claimed result}
\paragraph{Result 1.} For independent nearest-neighbor bond percolation on \(\mathbb Z^3\), write \(p_c^{\mathrm{bond}}\) for the first parameter permitting an infinite cluster. The source asserts
\[
 \mathbb P_{p_c^{\mathrm{bond}}}(\exists\text{ infinite open cluster})=0.
\]
For independent nearest-neighbor site percolation, it separately claims
\[
 \mathbb P_{p_c^{\mathrm{site}}}(\exists\text{ infinite open cluster})=0.
\]
These are two cubic-lattice statements at their respective critical parameters.

\paragraph{Result 2.} Suppose \(G\) is infinite, connected, locally finite, undirected, and has finitely many vertex orbits under automorphisms. For Bernoulli bond percolation, define \(p_c(G)\) as the infimum of parameters with positive probability of an infinite cluster. The claimed implication is
\[
 p_c(G)<1\quad\Longrightarrow\quad
 \mathbb P_{p_c(G)}(\exists\text{ infinite open cluster})=0.
\]
Equivalently, every fixed vertex \(x\) satisfies \(\mathbb P_{p_c}(|C_x|=\infty)=0\). The hypothesis \(p_c<1\) is essential to the stated scope.

\subsection{Outline of the argument}
Independent hyperedges provide a finite connection inequality that accommodates the two models. Finite-scale extension bounds allow an adaptive exploration to maintain successful connections; assuming critical percolation would make this construction survive below the threshold, a contradiction.

The proof separates growth regimes and uses finitary structural reduction in the polynomial-growth case. Two coordinate moves are glued through fresh corridors using a conditional connection inequality; adaptive exploration would then produce subcritical percolation if an infinite critical cluster existed.

\subsection{Scope relative to this problem}
For the lattice Z3, whose Bernoulli uniqueness threshold equals its critical threshold by amenable-lattice uniqueness, The cited theorem shows absence of a cluster at that endpoint; it supplies one negative case, not necessary-and-sufficient conditions for all quasi-transitive graphs. The cited theorem rules out bond clusters at pc<1; for amenable quasi-transitive graphs with pc=pu by the usual amenable uniqueness theorem this gives negative endpoint cases, but no full site/bond structural characterization at pu. This entry records partial support; complete coverage of the question is not established by this catalogue.

\subsection{Formal verification}
For family 213, the Lean scope covers cubic-lattice bond and site critical nonpercolation, and quasi-transitive bond graphs when $p_c<1$. The listed Comparator files are \texttt{CriticalPercolation.lean}, \texttt{CriticalZ3.lean}. Named declarations include \texttt{OAI.CriticalZ3.critical\_no\_infinite}. This site has not executed the Lean code. The scope descriptions do not certify the remaining manuscript arguments.

\subsection{Status caveat}
OpenAI attributes these manuscripts to work by an internal model. Its release cautions that arguments without formalization can contain mistakes. This catalogue has not independently checked the proofs; the labels report the mathematical claims attributed to OpenAI.

\subsection{Sources}
\begin{itemize}
\item OpenAI, \emph{Critical bond and site percolation on the cubic lattice}, 2026-09-24; citation key \texttt{OAI:Critical-bond-and-site-percolation-on-the-cubic-lattice-September-24-2026}. \href{https://github.com/openai/math/blob/adc7f1241b42e322a6451854ab7e4b4c146bf78a/preprints/Critical-bond-and-site-percolation-on-the-cubic-lattice-September-24-2026/paper.pdf}{PDF (GitHub)}; \href{https://github.com/openai/math/raw/adc7f1241b42e322a6451854ab7e4b4c146bf78a/preprints/Critical-bond-and-site-percolation-on-the-cubic-lattice-September-24-2026/paper.pdf}{PDF (direct)}; \href{https://github.com/openai/math/blob/adc7f1241b42e322a6451854ab7e4b4c146bf78a/preprints/Critical-bond-and-site-percolation-on-the-cubic-lattice-September-24-2026/README.md}{Manuscript page and citation}; \href{https://github.com/openai/math/blob/main/preprints/Critical-bond-and-site-percolation-on-the-cubic-lattice-September-24-2026/paper.pdf}{Latest version}; Statement source: \href{https://github.com/openai/math/blob/adc7f1241b42e322a6451854ab7e4b4c146bf78a/preprints/Critical-bond-and-site-percolation-on-the-cubic-lattice-September-24-2026/build/introduction.tex\#L17}{thm:main}.
\item OpenAI, \emph{No percolation at criticality on quasi-transitive graphs}, 2026-09-24; citation key \texttt{OAI:No-percolation-at-criticality-on-quasi-transitive-graphs-September-24-2026}. \href{https://github.com/openai/math/blob/adc7f1241b42e322a6451854ab7e4b4c146bf78a/preprints/No-percolation-at-criticality-on-quasi-transitive-graphs-September-24-2026/paper.pdf}{PDF (GitHub)}; \href{https://github.com/openai/math/raw/adc7f1241b42e322a6451854ab7e4b4c146bf78a/preprints/No-percolation-at-criticality-on-quasi-transitive-graphs-September-24-2026/paper.pdf}{PDF (direct)}; \href{https://github.com/openai/math/blob/adc7f1241b42e322a6451854ab7e4b4c146bf78a/preprints/No-percolation-at-criticality-on-quasi-transitive-graphs-September-24-2026/README.md}{Manuscript page and citation}; \href{https://github.com/openai/math/blob/main/preprints/No-percolation-at-criticality-on-quasi-transitive-graphs-September-24-2026/paper.pdf}{Latest version}; Statement source: \href{https://github.com/openai/math/blob/adc7f1241b42e322a6451854ab7e4b4c146bf78a/preprints/No-percolation-at-criticality-on-quasi-transitive-graphs-September-24-2026/build/sections/introduction.tex\#L26}{thm:main}.
\item Family 213: \href{https://github.com/openai/math/blob/adc7f1241b42e322a6451854ab7e4b4c146bf78a/lean/docs/213.md}{Lean scope}.
\item \href{https://github.com/openai/math/blob/adc7f1241b42e322a6451854ab7e4b4c146bf78a/lean/ComparatorChallenges/CriticalPercolation.json}{Comparator CriticalPercolation.json}.
\item \href{https://github.com/openai/math/blob/adc7f1241b42e322a6451854ab7e4b4c146bf78a/lean/ComparatorChallenges/CriticalPercolation.lean}{Comparator CriticalPercolation.lean}.
\item \href{https://github.com/openai/math/blob/adc7f1241b42e322a6451854ab7e4b4c146bf78a/lean/ComparatorChallenges/CriticalZ3.json}{Comparator CriticalZ3.json}.
\item \href{https://github.com/openai/math/blob/adc7f1241b42e322a6451854ab7e4b4c146bf78a/lean/ComparatorChallenges/CriticalZ3.lean}{Comparator CriticalZ3.lean}.
\item \href{https://github.com/openai/math/blob/adc7f1241b42e322a6451854ab7e4b4c146bf78a/CONTENTS.md}{OpenAI catalogue entry}.
\end{itemize}
\end{document}
